% answers_short.tex —— 2015 年·上海卷·理科 · 简短答案速查
% 由 tools/make_short_answers.py 自 answers.tex 的 \gksolhead/\gkanswer 提取，
% 并与 2015/2015_数学真题_上海理_答案要点.md 逐题比对后生成，勿手改（改源 MD 后重跑）。

\begin{gknotice}
\textbf{使用说明}：本节只给最终答案便于快速核对，完整推导见解析版。答案为编写者独立推导，\textbf{非官方标准答案}。
\end{gknotice}

\gkpart{填空题}
\begin{itemize}[leftmargin=2.6em]
  \item[\textbf{1.}] \(\{1,4\}\)
  \item[\textbf{2.}] \(\dfrac{1}{4}+\dfrac{1}{2}\mathrm{i}\)
  \item[\textbf{3.}] \(16\)
  \item[\textbf{4.}] \(4\)
  \item[\textbf{5.}] \(2\)
  \item[\textbf{6.}] \(\dfrac{\pi}{3}\)
  \item[\textbf{7.}] \(x=2\)
  \item[\textbf{8.}] \(120\)
  \item[\textbf{9.}] \(y=\pm\dfrac{\sqrt{3}}{2}x\)
  \item[\textbf{10.}] \(4\)
  \item[\textbf{11.}] \(45\)
  \item[\textbf{12.}] \(0.2\)
  \item[\textbf{13.}] \(8\)
  \item[\textbf{14.}] \(-\dfrac{16}{15}\)
\end{itemize}

\gkpart{选择题}
\begin{center}\begin{tabular}{ccccc}
\toprule
题号 & 15 & 16 & 17 & 18 \\
\midrule
答案 & B & D & B & A \\
\bottomrule
\end{tabular}\end{center}

\gkpart{解答题}
\begin{itemize}[leftmargin=2.6em]
  \item[\textbf{19.}] \(A_{1}\)，\(C_{1}\)，\(F\)，\(E\) 四点共面（它们所在平面为 \(x-y+z=1\)）；直线 \(CD_{1}\) 与平面 \(A_{1}C_{1}FE\) 所成角的正弦值为 \(\dfrac{\sqrt{15}}{15}\)。
  \item[\textbf{20.}] (1) \(t_{1}=\dfrac{3}{8}\) 小时；\(f(t_{1})=\dfrac{3\sqrt{41}}{8}\) 千米；(2) $f(t)=\begin{cases}\sqrt{25t^{2}-42t+18}, & \dfrac{3}{8}\leqslant t\leqslant\dfrac{7}{8},\\ 5-5t, & \dfrac{7}{8}\leqslant t\leqslant1.\end{cases}$；\(f(t)\) 在 \([t_{1},1]\) 上的最大值为 \(\dfrac{3\sqrt{41}}{8}\)（约 \(2.40\)，在 \(t=t_{1}=\dfrac{3}{8}\) 处取得），小于 \(3\)，故不超过 3
  \item[\textbf{21.}] (1) \(d=\dfrac{|x_{1}y_{2}-x_{2}y_{1}|}{\sqrt{x_{1}^{2}+y_{1}^{2}}}\)；且 \(S=2|x_{1}y_{2}-x_{2}y_{1}|\)；(2) \(S=\sqrt{2}\)
  \item[\textbf{22.}] (1) \(a_{n}=6n-5\)（\(n\in\mathbf{N}^{\ast}\)）；(2) 成立（\(b_{n}\leqslant b_{n_{0}}\) 对一切 \(n\in\mathbf{N}^{\ast}\)）；(3) \(\lambda\in\left(-\dfrac{1}{2},0\right)\)
  \item[\textbf{23.}] (1) 成立（\(h(x+6\pi)=h(x)+6\pi\)，故 \(\cos h(x+6\pi)=\cos h(x)\)，即 \(h\) 以 \(6\pi\) 为余弦周期）；(2) 成立（对任意 \(c\in[f(a),f(b)]\)，存在 \(x_{0}\in[a,b]\) 使 \(f(x_{0})=c\)；由单调性唯一）；(3) ① 两命题互为充要条件（\(u_{0}\in[0,T]\) 与 \(u_{0}+T\in[T,2T]\) 由 \(\cos f(u_{0}+T)=\cos f(u_{0})\) 逐点等价）；② 对任意 \(x\in[0,T]\) 有 \(f(x+T)=f(x)+f(T)=f(x)+4\pi\)
\end{itemize}
